\chapter{Choix du framework}
\label{chap:choix_framework}

Ce chapitre justifie le choix du \emph{framework} d'orchestration d'agents retenu pour l'implémentation. La démarche s'appuie sur le tableau comparatif détaillé conservé dans le dossier \texttt{Etudes/} (\texttt{comparaison\_frameworks\_IA\_detaille.xlsx}). La section~\ref{sec:criteres} formalise les critères d'évaluation, la section~\ref{sec:candidats} présente les candidats et leurs caractéristiques, la section~\ref{sec:decision} formule la décision motivée, et la section~\ref{sec:tradeoffs} explicite les compromis assumés.

\section{Critères d'évaluation}
\label{sec:criteres}

Les critères ont été construits à partir des concepts opératoires identifiés dans le glossaire technique du dossier d'études, en les hiérarchisant selon le besoin du cas d'étude.

\paragraph{Architecture en graphe orienté.}
Un graphe orienté est la structure la plus expressive pour modéliser un SMA : il permet d'identifier précisément qui communique avec qui, dans quel ordre, sous quelles conditions ; il est traçable, débogable, et proche de la théorie formelle des SMA. C'est un critère décisif pour un cas d'étude où l'auditabilité est une exigence.

\paragraph{État persistant typé (\emph{graph state}).}
Sans état persistant, un SMA oublie tout entre exécutions. Avec un état persistant, le système conserve une mémoire cohérente, peut être interrompu et repris (HITL, crash) et peut être audité \emph{post mortem}. La persistance doit être typée pour garantir l'intégrité des contrats inter-agents.

\paragraph{Paradigme de coordination explicite.}
Une coordination mal définie produit des agents redondants, contradictoires ou bloqués. Quatre paradigmes sont possibles : centralisé, hiérarchique, pair-à-pair, ou \emph{LLM-first} (le modèle décide en temps réel). Pour un cas d'étude critique, un modèle explicite et auditable est préférable à un modèle entièrement \emph{LLM-first}.

\paragraph{Human-in-the-loop natif.}
Mécanisme permettant à un humain d'intervenir en cours d'exécution. Crucial pour des applications où une mauvaise décision a des conséquences réelles, comme le contact d'un candidat. Le \emph{framework} doit supporter ce mécanisme nativement, sans contournement.

\paragraph{Intégration RAG et bases vectorielles.}
Le cas d'étude exige une mémoire vectorielle (calibration par fiche). Le \emph{framework} retenu doit s'interfacer simplement avec une base de type Chroma, Qdrant ou Pinecone.

\paragraph{Indépendance vis-à-vis du fournisseur.}
Le \emph{vendor lock-in} (dépendance à OpenAI, AWS, Microsoft) est un risque stratégique. Si le fournisseur augmente ses prix, arrête un service ou change ses conditions, l'infrastructure est à refaire. Pour un projet académique souverain, l'indépendance est un critère fort.

\paragraph{Licence et coût d'exploitation.}
Une licence permissive (MIT ou Apache 2.0) permet une utilisation libre. Le coût d'exploitation effectif dépend également des services associés : un \emph{framework} libre couplé à une plateforme propriétaire reproduit le \emph{lock-in}.

\paragraph{Courbe d'apprentissage.}
Pour un TER de quelques mois, une courbe d'apprentissage maîtrisable est nécessaire. Cependant, ce critère ne doit pas l'emporter sur la puissance d'expression : un \emph{framework} trop simple amputerait l'argumentation architecturale.

\paragraph{Maturité de l'écosystème.}
Documentation, communauté, intégrations, exemples, retours de production. Un \emph{framework} en \emph{alpha} est risqué pour un projet à délai contraint.

\section{Frameworks candidats}
\label{sec:candidats}

Neuf \emph{frameworks} sont confrontés aux critères. Les caractéristiques rapportées s'appuient sur le tableau comparatif détaillé du dossier d'études \cite{vaidhyanathan_taibi}.

\subsection{LangGraph}

Bibliothèque de la suite LangChain spécialisée dans la composition d'agents \cite{langgraph}. Caractéristiques : architecture en graphe orienté de premier ordre (nœuds, arêtes, branches conditionnelles, boucles) ; état persistant typé via \texttt{TypedDict} ; \emph{checkpointer} natif (en mémoire ou sur SQLite) ; mécanisme \texttt{interrupt\_before} dédié au HITL ; intégrations natives avec ChromaDB et autres bases vectorielles ; modèle d'orchestration explicite, indépendant du fournisseur LLM. Licence MIT. Courbe d'apprentissage moyenne à élevée.

\subsection{LangChain}

Couche de plus haut niveau \cite{langchain} pour construire des chaînes et agents. Très large adoption, écosystème riche. La structure d'une chaîne reste cependant linéaire ou faiblement branchante ; LangChain n'expose pas les mêmes primitives de graphe ni le même contrôle d'état que LangGraph. Souvent utilisée comme socle, complétée par LangGraph pour des topologies non triviales.

\subsection{AutoGen}

\emph{Framework} Microsoft Research \cite{autogen} centré sur la conversation multi-agents : les agents échangent des messages, un \emph{group chat} orchestre les tours. Supporte le HITL nativement. Très expressif pour des dialogues d'agents, mais le modèle conversationnel ne se prête pas naturellement à un pipeline en pas successifs avec étapes critiques (déduplication, scoring, routage par seuil).

\subsection{CrewAI}

\emph{Framework} récent \cite{crewai} centré sur la composition de rôles (\emph{crew}) et de tâches. S'appuie sur LiteLLM, donc indépendant du fournisseur. Faible courbe d'apprentissage. La coordination repose sur une orchestration souple (\emph{LLM-first}) qui rend l'audit pas-à-pas plus difficile que dans un modèle de graphe explicite.

\subsection{LlamaIndex}

\emph{Framework} \cite{llamaindex} excellent en RAG et indexation de corpus documentaires, avec une couche d'agents plus modeste. C'est davantage un moteur de recherche intelligent qu'un orchestrateur multi-agents.

\subsection{MetaGPT}

\emph{Framework} structuré autour de \emph{Standard Operating Procedures} (SOP) : enchaînement d'étapes prédéfinies imitant un cycle de production logicielle. Garantit la reproductibilité dans des environnements stables, mais la rigidité des SOP est mal adaptée à un pipeline de recrutement où le routage dépend dynamiquement des scores produits.

\subsection{OpenAI Swarm}

Bibliothèque très légère orientée \emph{handoffs} (passages de relais) entre agents. Simplicité didactique, mais : pas d'état partagé riche, pas de retour en arrière, dépendance forte à OpenAI (\emph{vendor lock-in}). Inadéquat pour le cahier des charges.

\subsection{AWS Strands}

\emph{Framework} centré sur AWS Bedrock, exploitant le protocole MCP. Adoption hors AWS limitée. Crée un \emph{vendor lock-in} sur AWS qui contredit l'objectif d'indépendance et de souveraineté.

\subsection{Semantic Kernel}

\emph{Framework} Microsoft fortement intégré à Azure et au stack .NET. Excellent dans son écosystème, mais reproduit la dépendance fournisseur sur Azure. Inadéquat pour un déploiement local et indépendant.

\section{Comparaison synthétique}

\begin{table}[H]
\centering
\footnotesize
\setlength{\tabcolsep}{4pt}
\begin{tabularx}{\textwidth}{@{}lcccccX@{}}
\toprule
\textbf{Framework} & \textbf{Graphe} & \textbf{État typé} & \textbf{HITL natif} & \textbf{Indépendance} & \textbf{Maturité} & \textbf{Adéquation au cas d'étude} \\
\midrule
LangGraph         & oui & oui & oui            & oui & élevée   & adéquat (architecture en graphe + HITL + RAG) \\
LangChain         & non & limité & partiel    & oui & élevée   & insuffisant pour topologie complexe \\
AutoGen           & non & non & oui            & oui & moyenne  & inadéquat : modèle conversationnel \\
CrewAI            & non & limité & partiel    & oui & moyenne  & inadéquat : audit pas-à-pas faible \\
LlamaIndex        & non & non & non            & oui & élevée   & complémentaire (RAG) mais pas orchestrateur \\
MetaGPT           & SOP & non & non            & oui & moyenne  & inadéquat : SOP trop rigides \\
OpenAI Swarm      & non & non & non            & non & faible   & inadéquat : \emph{lock-in} OpenAI \\
AWS Strands       & oui & oui & oui            & non & moyenne  & inadéquat : \emph{lock-in} AWS \\
Semantic Kernel   & non & oui & oui            & non & élevée   & inadéquat : \emph{lock-in} Azure \\
\bottomrule
\end{tabularx}
\caption{Comparatif synthétique des \emph{frameworks} confrontés aux critères du cas d'étude.}
\end{table}

\section{Décision motivée : LangGraph}
\label{sec:decision}

Le choix de LangGraph résulte de la conjonction de cinq éléments.

\paragraph{Adéquation structurelle.}
LangGraph est le seul \emph{framework} parmi les indépendants à offrir simultanément : une architecture en graphe orienté de premier ordre, un état persistant typé, un mécanisme HITL natif, une intégration directe à ChromaDB, et un \emph{checkpointer} permettant la reprise d'exécution. Tous les autres candidats indépendants (LangChain, CrewAI, AutoGen) excellent sur un sous-ensemble de critères mais n'offrent pas la combinaison complète.

\paragraph{Convergence avec l'argumentation théorique.}
Le chapitre~\ref{chap:etat_art} a montré que les SMA classiques apportent à l'Agentic AI des protocoles formels, des garanties de stabilité et une vérification explicite. Une architecture en graphe orienté est précisément le support naturel de ces apports : nœuds et arêtes formalisent les contrats, le typage protège les invariants, le \emph{checkpointer} permet l'audit. Le projet incarne ainsi l'apport SMA $\rightarrow$ Agentic AI dans son choix de \emph{framework}.

\paragraph{Indépendance et souveraineté.}
LangGraph n'est lié à aucun fournisseur LLM particulier : il s'interface aussi bien avec Ollama (local) qu'avec OpenAI, Anthropic ou tout fournisseur compatible LangChain. Cette indépendance est cohérente avec le choix d'Ollama local pour la confidentialité des données candidats.

\paragraph{Maturité et licence.}
Bibliothèque MIT, communauté active, documentation officielle riche, exemples nombreux. Ces éléments réduisent le risque temporel pour un TER à délai contraint.

\paragraph{Expressivité argumentative.}
Pour un jury évaluant l'argumentation architecturale, LangGraph fournit un vocabulaire précis et démontrable : \texttt{StateGraph}, \texttt{Send/Reduce}, \texttt{interrupt\_before}, \emph{conditional edges}, \emph{checkpointer}. Chacun de ces concepts cristallise un patron de coordination qui peut être exhibé et défendu.

\section{Tradeoffs assumés}
\label{sec:tradeoffs}

Trois compromis sont explicitement assumés par ce choix.

\paragraph{Courbe d'apprentissage.}
LangGraph est plus exigeant qu'un \emph{framework} \emph{LLM-first} comme CrewAI. Il faut comprendre les notions de \emph{StateGraph}, de \emph{reducer}, de \emph{checkpointer}, de \emph{conditional edges}. Cet investissement initial est compensé par la lisibilité du code obtenu : un nouveau lecteur peut reconstruire mentalement le pipeline en lisant le fichier \texttt{src/graph.py} (environ deux cents lignes).

\paragraph{Verbosité du typage d'état.}
Le \texttt{TypedDict} \texttt{GraphState} comporte une quinzaine de champs explicites. Ce typage explicite est plus verbeux qu'un état implicite à la AutoGen, mais chaque champ documente précisément le contrat entre agents et permet une détection statique des fautes (champ inexistant, type incompatible).

\paragraph{Dépendance à l'écosystème LangChain.}
LangGraph s'inscrit dans la suite LangChain, dont l'évolution rapide implique des ruptures occasionnelles d'API. Le projet épingle des versions précises (\texttt{langgraph>=1.1.0}, \texttt{langchain>=0.3.0}) et délimite l'usage de LangChain au strict nécessaire (initialisation du modèle via \texttt{init\_chat\_model}). Le couplage est suffisamment limité pour qu'une migration ultérieure soit envisageable.

\medskip

En conclusion, LangGraph satisfait simultanément les exigences d'expressivité (graphe, état, routage), d'auditabilité (typage, \emph{checkpointer}), de gouvernance (HITL natif), d'indépendance (pas de \emph{vendor lock-in}) et de cohérence avec l'argumentation théorique du chapitre précédent. Le chapitre suivant détaille comment ce choix se concrétise dans le cas d'étude.
